\section{Compiling the renderer}
\label{sec:compiling}
Mitsuba is built with CMake 3.16 or later and a C++11-compatible compiler.
The following steps describe the common build process on Linux, Windows, and macOS.

\subsection{Common steps}
To get started, download a recent version of the Mitsuba source code. Before
doing this, ensure that you have read the licensing agreement
(\secref{license}), and that you abide by its contents. Note that, being a ``viral''
license, the GPL automatically applies to derivative work. Amongst other things, this
means that Mitsuba's source code is \emph{off-limits} to those who develop rendering
software not distributed under a compatible license.

Check that Git (\url{https://git-scm.com/}) and CMake are installed, then run:
\begin{shell}
$\texttt{\$}$ git clone https://github.com/mitsuba-renderer/mitsuba
$\texttt{\$}$ cd mitsuba
\end{shell}

\subsection{Dependencies}
The repository includes a precompiled \code{dependencies} bundle for some third-party
libraries on supported platforms. It is only a partial set: libraries not included in
that bundle must be installed separately or made available to CMake. In particular,
Boost 1.80 or later is required to build Mitsuba. On Windows, the bundle can be placed
in the source directory as \code{dependencies}; a different location can be selected
with \code{-DMITSUBA_DEPENDENCIES_ROOT=...}.

The required dependencies include OpenGL, Xerces-C, Eigen 3, GLEW, zlib, and Boost
1.80 or later. Linux builds additionally require X11 and Xxf86vm development files.
Package names vary by operating system and distribution.

The graphical interface is enabled by default when Qt is found. Building it requires
Qt 5.9 or later, including the Core, Gui, Widgets, OpenGL, Xml, XmlPatterns, and Network
modules. If Qt is unavailable, CMake skips the GUI; to disable it explicitly, pass
\code{-DMITSUBA_BUILD_GUI=OFF}.

PNG, JPEG, OpenEXR, and FFTW support is enabled when their development libraries are
available. The Python module is optional and requires Python 3 development files and
the matching Boost.Python library; enable it with \code{-DMITSUBA_BUILD_PYTHON=ON}.

\subsection{Building}
Configure and build from the root of the source tree:
\begin{shell}
$\texttt{\$}$ cmake -S . -B build -DCMAKE_BUILD_TYPE=Release
$\texttt{\$}$ cmake --build build --parallel
\end{shell}
To install the resulting binaries and libraries, run:
\begin{shell}
$\texttt{\$}$ cmake --install build --prefix <install-prefix>
\end{shell}
On Windows, run these commands from a Visual Studio developer command prompt, or
select the desired Visual Studio generator when configuring CMake.

\subsection{Compilation options}
\label{sec:compiling-flags}
The most commonly adjusted options are passed to CMake during configuration:
\begin{description}
\item[\code{MITSUBA_SPECTRUM_SAMPLES}] Number of spectral samples; defaults to 3.
\item[\code{MITSUBA_SINGLE_PRECISION}] Use single-precision calculations; enabled by default.
\item[\code{MITSUBA_ENABLE_SSE}] Enable SSE code paths; enabled by default.
\item[\code{MITSUBA_ENABLE_COHERENT_RT}] Enable coherent ray tracing; enabled by default.
\item[\code{MITSUBA_DEBUG_CHECKS}] Enable debug checks; enabled by default.
\item[\code{MITSUBA_BUILD_GUI}] Build the Qt graphical interface when Qt is available; enabled by default.
\item[\code{MITSUBA_BUILD_COLLADA}] Build COLLADA support when COLLADA DOM is available; enabled by default.
\item[\code{MITSUBA_BUILD_PYTHON}] Build the Python 3 module; disabled by default.
\end{description}